\chapter{The equatorial disk and critical catenoid}\label{sec:examples-free-boundary}
These examples connect the catenoid calculation in Example \ref{ex:explicit-catenoid} to the free-boundary geometry of Section \ref{sec:free-boundary-ball}. The critical catenoid and its Steklov spectrum are discussed in \cite{FS11}, Section 3, pp.~6--9. We derive the boundary condition and scale directly, and distinguish a negative variation from a calculation of the full Morse index.

\begin{example}[The equatorial ball]\label{ex:equatorial-free-boundary-ball}
The flat unit ball $D^n=\overline B^{n+1}\cap\{x_{n+1}=0\}$ is free-boundary minimal. It has volume $\omega_n$, boundary volume $n\omega_n$, and $A=0$. A constant normal speed has $Q(1,1)=-n\omega_n$. Each coordinate function $x_i$, $1\leq i\leq n$, is a Jacobi field for the free-boundary second variation, meaning $Q(x_i,f)=0$ for every smooth $f$.
\end{example}
\begin{proof}
Flatness gives minimality, and the outward conormal of the unit ball is its position vector. The volume formulas and Proposition \ref{prop:free-boundary-normal-form} give the negative variation. Since $\Delta x_i=0$ and $\partial_\eta x_i=x_i$, integration by parts gives
$Q(x_i,f)=\int_D\langle\nabla x_i,\nabla f\rangle-\int_{\partial D}x_if=0$.
These fields arise from rotations that tilt the equatorial plane; the constant speed does not arise from such a symmetry.
\end{proof}

\begin{proposition}[The Steklov spectrum of the flat disk]\label{prop:disk-steklov-spectrum}
For the unit disk $D^2$ with its flat metric, the Steklov eigenvalues are $0,1,2,\ldots$. The zero eigenspace is the constants; for integer $k\geq1$, the eigenvalue $k$ has basis $r^k\cos(k\theta)$ and $r^k\sin(k\theta)$.
\end{proposition}
\begin{proof}
A smooth harmonic function on the closed disk has the harmonic Fourier expansion
\[
 u(r,\theta)=a_0+\sum_{k\geq1}r^k(a_k\cos k\theta+b_k\sin k\theta).
\]
This follows by taking Fourier coefficients in the harmonic equation: the radial ODE has solutions $r^{\pm k}$, of which only $r^k$ is smooth at the origin; its zero mode is constant. The boundary condition $u_r=\sigma u$ at $r=1$ makes $(k-\sigma)a_k=(k-\sigma)b_k=0$ and $\sigma a_0=0$. Fourier uniqueness proves the assertion, while each displayed function directly realizes its eigenvalue.
\end{proof}

\begin{proposition}[Index one and nullity two for the equatorial disk]\label{prop:equatorial-disk-index}
The equatorial unit disk in $B^3$ has free-boundary Morse index one and nullity two. Here the index is the largest dimension of a subspace on which $Q$ is negative definite, and nullity is the dimension of the bilinear kernel $\{f:Q(f,g)=0\text{ for every }g\}$.
\end{proposition}
\begin{proof}
If the boundary trace of $f$ is $a_0+\sum_{k\geq1}(a_k\cos k\theta+b_k\sin k\theta)$, let $h$ be its harmonic extension. Green's identity gives
$\int_D|\nabla f|^2=\int_D|\nabla h|^2+\int_D|\nabla(f-h)|^2$.
Orthogonality of the trigonometric modes yields
\[
 Q(f,f)=\int_D|\nabla(f-h)|^2-2\pi a_0^2
       +\pi\sum_{k\geq1}(k-1)(a_k^2+b_k^2).
\]
The codimension-one subspace of functions with zero boundary mean is nonnegative for $Q$, so a negative-definite subspace has dimension at most one. The constant function gives a negative direction, proving index one. If $Q(f,g)=0$ for all $g$, interior tests imply $\Delta f=0$ and the boundary terms give $\partial_\eta f=f$. Proposition \ref{prop:disk-steklov-spectrum} identifies its solutions as the span of $x_1,x_2$. Conversely those functions are in the kernel by Example \ref{ex:equatorial-free-boundary-ball}. The zero cone $Q(f,f)=0$ is not being identified with this bilinear kernel.
\end{proof}

\begin{lemma}[The critical catenoid parameter]\label{lem:critical-catenoid-parameter}
There is a unique $t_0>0$ such that $t_0\tanh t_0=1$, and $t_0>1$. Put
\[
 a=\frac1{t_0\cosh t_0}.
\]
Then $a=(\cosh^2t_0+t_0^2)^{-1/2}$, $a\cosh t_0=1/t_0$, and $at_0=\operatorname{sech}t_0$.
\end{lemma}
\begin{proof}
The function $t\tanh t$ starts at zero, has derivative $\tanh t+t\operatorname{sech}^2t>0$ for $t>0$, and tends to infinity. Its value at $1$ is less than one. The root exists uniquely and exceeds one. From $\tanh t_0=1/t_0$ and $\operatorname{sech}^2t_0=1-1/t_0^2$, one gets $\cosh^2t_0+t_0^2=t_0^2\cosh^2t_0$, proving the scale identities. This is the centered symmetric case of \cite{FS11}, Theorem 3.1.
\end{proof}

\begin{example}[A free-boundary catenoidal annulus]\label{ex:critical-catenoid}
With $a,t_0$ as above, the immersion
\[
 F(t,\theta)=a(\cosh t\cos\theta,\cosh t\sin\theta,t),
 \qquad -t_0\leq t\leq t_0,\quad \theta\in\mathbb R/(2\pi\mathbb Z),
\]
is a smooth embedded free-boundary minimal annulus in the unit ball. Its boundary circles have radius $1/t_0$ and heights $\pm\operatorname{sech}t_0$. Its metric is $a^2\cosh^2t\,(dt^2+d\theta^2)$.
\end{example}
\begin{proof}
Example \ref{ex:explicit-catenoid} gives the metric, minimality, and embeddedness. Since $|F|^2=a^2(\cosh^2t+t^2)$ is strictly increasing in $|t|>0$, the scale identity gives $|F|=1$ exactly at $t=\pm t_0$ and $|F|<1$ in the interior. At the upper circle the unit conormal is
\[
 \frac{F_t}{|F_t|}=(\tanh t_0\cos\theta,\tanh t_0\sin\theta,\operatorname{sech}t_0)=F(t_0,\theta).
\]
At the lower circle the outward conormal is $-F_t/|F_t|$, which similarly equals $F(-t_0,\theta)$. The scale identities give the stated circle radii and heights. This establishes the free-boundary condition directly.
\end{proof}

\begin{proposition}[Area, boundary length, and an unstable direction]\label{prop:critical-catenoid-integrals}
For the critical catenoid $C$,
\[
 |C|=\frac{2\pi}{t_0},\qquad |\partial C|=\frac{4\pi}{t_0},\qquad
 \int_C|A|^2=\frac{8\pi}{t_0},\qquad Q(1,1)=-\frac{12\pi}{t_0}.
\]
These computations exhibit a negative variation but do not determine the full Morse index.
\end{proposition}
\begin{proof}
The area integral is
$2\pi a^2\int_{-t_0}^{t_0}\cosh^2t\,dt=2\pi a^2(t_0+\sinh t_0\cosh t_0)=2\pi/t_0$.
For the simplification, use $a^2=1/(t_0^2\cosh^2t_0)$ and $\sinh t_0=\cosh t_0/t_0$. The two boundary circles give length $4\pi/t_0$, also checking $|\partial C|=2|C|$. By Proposition \ref{prop:catenoid-curvature-example}, $|A|^2=2/(a^2\cosh^4t)$. Multiplication by the area element gives $2\operatorname{sech}^2t\,dt\,d\theta$, whose integral is $8\pi\tanh t_0=8\pi/t_0$. Substitute into the normal quadratic form.
\end{proof}

\begin{proposition}[First Steklov eigenspace of the critical catenoid]\label{prop:critical-catenoid-steklov}
The critical catenoid has first positive Steklov eigenvalue $1$, of multiplicity three, with eigenspace spanned by its three ambient coordinate functions. In particular $\sigma_1|\partial C|=4\pi/t_0$.
\end{proposition}
\begin{proof}
In dimension two, the conformal factor in the metric cancels from the harmonic equation. Thus it is $u_{tt}+u_{\theta\theta}=0$. At either boundary the unit outward derivative is $\pm t_0\partial_t$, since $a\cosh t_0=1/t_0$. The angular zero mode is affine in $t$; its constant solution has eigenvalue zero, and its odd solution $t$ has eigenvalue one. For integer $k\geq1$, the even modes $\cosh(kt)\cos(k\theta)$ and $\cosh(kt)\sin(k\theta)$ have eigenvalue
\[
 \sigma_k^{\rm even}=t_0k\tanh(kt_0),
\]
while the analogous odd modes with $\sinh(kt)$ have eigenvalue
$\sigma_k^{\rm odd}=t_0k\coth(kt_0)$.
Fourier decomposition and reflection into even and odd parts exhaust the smooth harmonic eigenfunctions. The function $s\mapsto s\tanh(st_0)$ is strictly increasing for $s>0$. Hence its smallest integer-mode value is $\sigma_1^{\rm even}=t_0\tanh t_0=1$; all $k\geq2$ even values exceed one. Every odd value is greater than $t_0k>1$. Together with the zero-angular-mode function $t$, this proves precisely three eigenfunctions at one, which are multiples of $F_1,F_2,F_3$. This specializes the separation of variables in \cite{FS11}, pp.~7--9; it makes no claim about the sharp bound among all annular metrics.
\end{proof}
